\clearpage
\VerdierPhysicalPageAnchor{0231}%
\begingroup\small\setlength{\parskip}{0pt}\setlength{\emergencystretch}{1.5em}
\setlength{\abovedisplayskip}{4pt plus 2pt minus 2pt}%
\setlength{\belowdisplayskip}{4pt plus 2pt minus 2pt}%

\VerdierSemanticAnchor{ch4d-4-1-7-proof}
まず、$X$ が二重複体で整数 $m=1$ の場合に命題が証明済みであると仮定し、
この場合へ帰着する方法を示す。$\varphi\colon[n]\longrightarrow[2]$ を写像

\VerdierSemanticAnchor{ch4d-4-1-7-1}
\begin{equation}
\left\{
\begin{aligned}
\varphi(m)&=1,\\
\varphi(i)&=2,\qquad i\in[n],\quad i\neq m.
\end{aligned}
\right.
\tag{4.1.7.1}
\end{equation}

とする。$\int_\varphi X=Z$（第 I 章、2.2.2）と置く。二重複体 $Z$ に
(4.1.6.$i$)、$1\leq i\leq4$ で述べた構成を施すことができる。二重複体

\VerdierSemanticAnchor{ch4d-4-1-7-2}
\begin{equation}
{}^{1}\!Z(p,q),\qquad
(p,q)\in\mathcal{F}\ell_1(\widetilde{\mathbb Z}),
\tag{4.1.7.2}
\end{equation}

および二重複体の射

\VerdierSemanticAnchor{ch4d-4-1-7-3}
\begin{equation}
\begin{aligned}
f_Z\bigl((p,q),(p',q')\bigr)&\colon
{}^{1}\!Z(p,q)\longrightarrow{}^{1}\!Z(p',q'),\\
g_Z(p,q,r)&\colon
{}^{1}\!Z(q,r)\longrightarrow T_1\bigl({}^{1}\!Z(p,q)\bigr).
\end{aligned}
\tag{4.1.7.3}
\end{equation}

が得られる。このとき

\VerdierSemanticAnchor{ch4d-4-1-7-4}
\begin{equation}
\left\{
\begin{aligned}
\int_\varphi{}^{m}\!X(p,q)
  &={}^{1}\!Z(p,q),
  &&(p,q)\in\mathcal{F}\ell_1(\widetilde{\mathbb Z}),\\
\int_\varphi f\bigl((p,q),(p',q')\bigr)
  &=f_Z\bigl((p,q),(p',q')\bigr),
\end{aligned}
\right.
\tag{4.1.7.4}
\end{equation}

であり（直接確かめられる）、さらに射 $g_Z(p,q,r)$ は合成射

\VerdierSemanticAnchor{ch4d-4-1-7-5}
\begin{equation}
\int_\varphi{}^{m}\!X(q,r)
\xrightarrow{\int_\varphi g(p,q,r)}
\int_\varphi T_m\bigl({}^{m}\!X(p,q)\bigr)
\xrightarrow{(**)}
T_1\left(\int_\varphi{}^{m}\!X(p,q)\right)
\tag{4.1.7.5}
\end{equation}

である。ここで $(**)$ は関手 $\int_\varphi$ とシフトとの可換同型
（第 I 章、2.3.2.1）である（これを確かめるには、$Z$ の微分の明示的な定義
（第 I 章、2.2.2.5）と可換同型の定義（第 I 章、2.3.2.1）を用いる）。
次に、第一次数でフィルター付けされた $Z$ によって定まるスペクトル対象を
構成する。等式 $\int_2\int_\varphi=\int_n$、(4.1.7.4)、(4.1.7.5)、および
（第 I 章、2.3.3）から、
$\bigl({}^{m}\!Y(p,q),{}^{m}\!d(p,q,r)\bigr)$ はこのスペクトル対象に等しい。
\endgroup
